\section{Conclusion}\label{sec:conclusion}
Cable geometry makes braking authority a controlled state: the team brakes only as hard as the cable directions allow, and these change at a bounded rate. An initially feasible HOCBF can lose feasibility when the geometry cannot change quickly enough, under the stated gain-dependent condition. A support-function calculation and an explicit swing-and-brake maneuver give useful continuous-time inner and outer viability bounds, and the argument extends to every system with second-order authority. The inner bound is conservative where the cables must first swing: safe trajectories were found from as little as $0.25\,\Dres$ there, against $0.97$--$0.98\,\Dres$ with the cables already behind the payload; the reserved swing acceleration accounts for part of this gap.

The simulations also reveal that the current sampled filter and physical cable model do not satisfy the theorem's hypotheses: the barrier lost up to $0.85\,\mathrm m$ within one $5\,\mathrm{ms}$ hold; of 100 trials started near the boundary of the safe set, 56 had a slack cable and 59 reached the wall. A certified sampled-data bound, coordinated force allocation, and reproducible full-order validation are the next requirements for a deployment claim.

